  \hspace{3mm}

\centerline{\bf \Huge ТЧ и алгебра с Олимпиады им. Эйлера}

\hspace{3mm}

\problem{1} Даны два ненулевых числа. Если к каждому из них прибавить единицу, а также из каждого из них вычесть единицу, то сумма обратных величин четырёх полученных чисел будет равна 0. Какое число
может получиться, если из суммы исходных чисел вычесть сумму
их обратных величин? Найдите все возможности.

\problem{2} Числа $x$ и $y$, не равные $0$, удовлетворяют неравенствам $x^2 - x > y^2$ и $y^2 - y  > x^2$. Какой знак может иметь произведение $xy$ (укажите все возможности)?

\problem{3} Будем называть натуральное число красивым, если в его десятичной записи поровну цифр 0, 1, 2, а других цифр нет (во избежание недоразумений напомним, что десятичная запись числа не может начинаться с нуля). Может ли произведение двух красивых чисел быть красивым?

\problem{4}  Докажите, что существует натуральное число n, большее $10^{100}$, такое, что сумма всех
простых чисел, меньших n, взаимно проста с n.

\problem{5} Пусть $p_1,\ p_2, \dotsc , p_{99},\ p_{100} -$ сто простых чисел, среди которых нет одинаковых. Натуральные числа $a_1, a_2, \dotsc, a_k$ большие 1, таковы, что каждое из чисел $p_1p_2^3, \ p_2p_3^3, \dotsc,\\ p_{99}p_{100}^3,\ p_{100}p_1^3$ равно произведению каких-то двух из
чисел $a_1, a_2, \dotsc, a_k$. Докажите, что $k\geq 150.$

\problem{6} Можно ли число 240 представить в виде суммы девяти двузначных чисел (среди которых могут быть и одинаковые), в десятичной записи каждого из которых есть девятка?